\documentclass[11pt]{article}

% --- Packages ---
\usepackage[utf8]{inputenc}
\usepackage[T1]{fontenc}
\usepackage{geometry}
\geometry{a4paper, margin=2.5cm}
\usepackage{booktabs}
\usepackage{longtable}
\usepackage{array}
\usepackage{caption}
\usepackage{hyperref}

% Column types for wrapping
\newcolumntype{L}[1]{>{\raggedright\arraybackslash}p{#1}}
\newcolumntype{C}[1]{>{\centering\arraybackslash}p{#1}}

\title{Supplementary Materials\\[0.5em]
\large A Quality of Care Index for Drug-Susceptible Tuberculosis:\\
A Systematic Analysis Using the Global Burden of Disease Study 2021}
\date{}

\begin{document}

\maketitle

% ====================================================================
%  TABLE S1: GATHER Checklist
% ====================================================================
\setcounter{table}{0}
\renewcommand{\thetable}{S\arabic{table}}

\begin{longtable}{C{1.2cm} L{4.8cm} L{8.5cm}}
\caption{GATHER Checklist: Compliance of the present study with the Guidelines for Accurate and Transparent Health Estimates Reporting.}
\label{tab:gather} \\
\toprule
\textbf{Item} & \textbf{GATHER Recommendation} & \textbf{Compliance in This Study} \\
\midrule
\endfirsthead

\multicolumn{3}{l}{\textit{Table~S1 continued from previous page}} \\
\toprule
\textbf{Item} & \textbf{GATHER Recommendation} & \textbf{Compliance in This Study} \\
\midrule
\endhead

\midrule
\multicolumn{3}{r}{\textit{Continued on next page}} \\
\endfoot

\bottomrule
\endlastfoot

1 &
Define the indicator(s), populations, time period(s), and geography(ies) covered. &
The Quality of Care Index (QCI) is defined as a composite indicator for drug-susceptible tuberculosis care quality across 194 countries and territories, 31 Iranian provinces, 5 age groups, both sexes, and the period 1990--2021. \\
\addlinespace

2 &
List the funding sources for the work. &
Funding sources are stated in the main manuscript acknowledgements section. \\
\addlinespace

3 &
Describe how the data were identified and accessed. &
All input data were obtained from the Global Burden of Disease (GBD) 2021 study via the IHME Global Health Data Exchange (\url{https://ghdx.healthdata.org/}). Indicators retrieved include incidence, prevalence, mortality, YLLs, YLDs, and DALYs for drug-susceptible tuberculosis. \\
\addlinespace

4 &
Specify the inclusion and exclusion criteria, including any geographical, temporal, or population restrictions. &
All 194 countries and territories available in GBD 2021 were included. Sub-national data were restricted to the 31 provinces of Iran. The study period spans 1990--2021. Age groups analysed were: ${<}5$, 5--14, 15--49, 50--69, 70+~years, age-standardised, and all ages. Both sexes were analysed separately and combined. \\
\addlinespace

5 &
Provide enough detail for the study to be replicated; describe and provide sources for all data inputs, analytical steps, and modelling methods. &
Four indicator ratios (MIR, YLL-to-YLD, DALY-to-prevalence, prevalence-to-incidence) were derived from GBD 2021 estimates. Principal Component Analysis (PCA) was trained on age-standardised, both-sexes data and applied to all demographic strata. The first principal component was rescaled to a 0--100 index using global min--max normalisation. Full methodological details, including the PCA training strategy and scaling approach, are provided in the Methods section and this supplement (Table~S2). \\
\addlinespace

6 &
Describe the methods used to quantify uncertainty of the estimates. &
Monte Carlo simulation with 1{,}000 draws was used to propagate GBD uncertainty through the PCA and rescaling pipeline. For each draw, input indicator ratios were sampled from their GBD-provided uncertainty distributions, PCA loadings from the base model were applied, and the resulting scores were rescaled. The 2.5th and 97.5th percentiles of the resulting distribution define the 95\% uncertainty intervals reported for all QCI estimates. \\
\addlinespace

7 &
Describe and provide results for each analysis step. &
Results are presented at global, regional (World Bank regions, SDI quintiles), national, and sub-national (Iranian provinces) levels. QCI values with 95\% uncertainty intervals are reported for 1990 and 2021. Temporal trends are characterised by the Average Annual Percentage Change (AAPC) with 95\% confidence intervals. \\
\addlinespace

8 &
Provide quantitative measures of the uncertainty of the estimates. &
All QCI estimates are accompanied by 95\% uncertainty intervals derived from Monte Carlo simulation. AAPC estimates include 95\% confidence intervals from log-linear regression. Sensitivity analyses comparing 3- versus 4-component PCA models are reported (Table~S2). \\
\addlinespace

9 &
Describe the limitations of the data sources and analytical methods. &
Limitations include the ecological nature of the analysis, reliance on GBD modelled estimates which may have data sparsity in low-income settings, the composite index masking heterogeneity among component indicators, and the cross-sectional SDI classification that does not capture within-period socio-demographic transitions. These are discussed in the Limitations section of the main manuscript. \\
\addlinespace

10 &
Provide data and code so that results can be reproduced. &
All input data are publicly available from the IHME GBD Results Tool. Analytical code and derived datasets will be made available upon publication in a public repository. \\

\end{longtable}


\clearpage

% ====================================================================
%  TABLE S2: Sensitivity Analysis Summary
% ====================================================================

\begin{longtable}{L{5.5cm} L{9cm}}
\caption{Sensitivity Analysis Summary: robustness of the Quality of Care Index to methodological decisions.}
\label{tab:sensitivity} \\
\toprule
\textbf{Analysis} & \textbf{Result} \\
\midrule
\endfirsthead

\multicolumn{2}{l}{\textit{Table~S2 continued from previous page}} \\
\toprule
\textbf{Analysis} & \textbf{Result} \\
\midrule
\endhead

\midrule
\multicolumn{2}{r}{\textit{Continued on next page}} \\
\endfoot

\bottomrule
\endlastfoot

\multicolumn{2}{l}{\textbf{A.\quad PCA component selection: 3 vs.\ 4 indicators}} \\
\addlinespace

PC1 variance explained (3-component model: MIR, YLL-to-YLD, DALY-to-prevalence) &
87.98\% of total variance captured by PC1. \\
\addlinespace

PC1 variance explained (4-component model: adds prevalence-to-incidence) &
66.13\% of total variance captured by PC1. \\
\addlinespace

Contribution of prevalence-to-incidence (PERtoINC) to PC1 &
Only 0.37\% contribution to the first principal component in the 4-component model, indicating near-zero discriminatory value. \\
\addlinespace

Rank correlation between 3- and 4-component QCI &
Spearman $\rho = 0.998$, confirming near-perfect rank agreement between the two models. \\
\addlinespace

Country ranking overlap (top 20) &
19 of 20 top-ranked countries were identical between the 3- and 4-component models. \\
\addlinespace

Country ranking overlap (bottom 20) &
20 of 20 bottom-ranked countries were identical between the two models. \\
\addlinespace

Decision &
The 3-component model (MIR, YLL-to-YLD, DALY-to-prevalence) was selected as the primary specification due to higher variance concentration in PC1 and the negligible contribution of PERtoINC. \\
\addlinespace
\midrule
\addlinespace

\multicolumn{2}{l}{\textbf{B.\quad PCA training strategy}} \\
\addlinespace

Training data &
PCA was trained on age-standardised, both-sexes data pooled across all 194 countries and all years (1990--2021). This single set of loadings was then applied to all demographic strata (age groups, sex-specific estimates) to ensure a consistent scoring framework. \\
\addlinespace

Rationale &
Training on age-standardised data avoids confounding by age structure and ensures that the principal component weights reflect quality-of-care variation rather than demographic composition. Applying a single set of loadings across strata allows direct comparability of QCI values across ages and sexes. \\
\addlinespace
\midrule
\addlinespace

\multicolumn{2}{l}{\textbf{C.\quad Scaling approach}} \\
\addlinespace

Normalisation method &
Global min--max scaling was applied to the raw PC1 scores, transforming them to a 0--100 scale where 100 represents the highest observed quality of care. The global minimum and maximum were computed across all countries, years, and demographic strata jointly. \\
\addlinespace

Rationale &
A single global reference range ensures that QCI values are directly comparable across time, geography, and population subgroups. Country-specific or year-specific rescaling would preclude temporal and cross-national comparisons. \\

\end{longtable}

\clearpage

% ====================================================================
%  TABLE S3 placeholder -- generated by Python script
% ====================================================================

% The following table is auto-generated by generate_supplementary_table.py
% Include it with: \begin{longtable}{r l r r r r l l}
\caption{Complete country ranking by Quality of Care Index in 2021 (age-standardised, both sexes), with QCI in 1990, absolute change, Average Annual Percentage Change (AAPC) over 1990--2021, and SDI group.}
\label{tab:all_countries} \\
\toprule
\textbf{Rank} & \textbf{Country} & \textbf{QCI 2021} & \textbf{QCI 1990} & \textbf{Change} & \textbf{AAPC} & \textbf{95\% CI} & \textbf{SDI Group} \\
\midrule
\endfirsthead

\multicolumn{8}{l}{\textit{Table~S3 continued from previous page}} \\
\toprule
\textbf{Rank} & \textbf{Country} & \textbf{QCI 2021} & \textbf{QCI 1990} & \textbf{Change} & \textbf{AAPC} & \textbf{95\% CI} & \textbf{SDI Group} \\
\midrule
\endhead

\midrule
\multicolumn{8}{r}{\textit{Continued on next page}} \\
\endfoot

\bottomrule
\endlastfoot
1 & Bermuda & 99.94 & 99.65 & 0.29 & 0.01 & (0.01, 0.01) & High \\
2 & New Zealand & 99.63 & 97.38 & 2.25 & 0.09 & (0.07, 0.10) & High \\
3 & Malta & 99.54 & 98.60 & 0.94 & 0.04 & (0.03, 0.04) & High \\
4 & Antigua and Barbuda & 99.31 & 98.02 & 1.28 & 0.04 & (0.04, 0.05) & High-middle \\
5 & Luxembourg & 99.16 & 97.62 & 1.54 & 0.06 & (0.06, 0.07) & High \\
6 & Cuba & 99.04 & 97.89 & 1.15 & 0.04 & (0.03, 0.05) & High-middle \\
7 & Taiwan & 99.01 & 93.34 & 5.67 & 0.13 & (0.11, 0.16) & High \\
8 & Sweden & 98.99 & 95.87 & 3.12 & 0.13 & (0.12, 0.14) & High \\
9 & Switzerland & 98.94 & 96.68 & 2.26 & 0.09 & (0.08, 0.09) & High \\
10 & Grenada & 98.91 & 96.49 & 2.42 & 0.07 & (0.06, 0.08) & Middle \\
11 & Czech Republic & 98.79 & 96.86 & 1.92 & 0.06 & (0.05, 0.06) & High \\
12 & Slovenia & 98.70 & 96.64 & 2.07 & 0.07 & (0.06, 0.08) & High \\
13 & Denmark & 98.70 & 97.09 & 1.62 & 0.06 & (0.06, 0.07) & High \\
14 & Australia & 98.70 & 97.36 & 1.34 & 0.05 & (0.05, 0.06) & High \\
15 & Italy & 98.69 & 96.69 & 1.99 & 0.07 & (0.05, 0.08) & High \\
16 & China & 98.68 & 94.21 & 4.47 & 0.14 & (0.12, 0.15) & High-middle \\
17 & Mauritius & 98.64 & 96.71 & 1.94 & 0.05 & (0.04, 0.06) & High-middle \\
18 & Singapore & 98.61 & 95.64 & 2.97 & 0.10 & (0.09, 0.10) & High \\
19 & Armenia & 98.60 & 97.09 & 1.51 & 0.09 & (0.07, 0.11) & Middle \\
20 & Belgium & 98.58 & 96.62 & 1.96 & 0.08 & (0.07, 0.08) & High \\
21 & Germany & 98.49 & 95.65 & 2.84 & 0.10 & (0.09, 0.11) & High \\
22 & Moldova & 98.45 & 97.52 & 0.92 & 0.06 & (0.04, 0.09) & Middle \\
23 & United Kingdom & 98.44 & 97.42 & 1.02 & 0.05 & (0.04, 0.05) & High \\
24 & Canada & 98.35 & 95.04 & 3.31 & 0.11 & (0.09, 0.13) & High \\
25 & Netherlands & 98.28 & 96.57 & 1.70 & 0.08 & (0.07, 0.09) & High \\
26 & Georgia & 98.26 & 95.34 & 2.92 & 0.12 & (0.10, 0.14) & Middle \\
27 & Lebanon & 98.24 & 95.90 & 2.34 & 0.08 & (0.07, 0.08) & Middle \\
28 & Austria & 98.19 & 95.66 & 2.52 & 0.08 & (0.07, 0.10) & High \\
29 & Barbados & 98.16 & 96.78 & 1.38 & 0.06 & (0.05, 0.07) & High \\
30 & Sri Lanka & 98.13 & 93.59 & 4.54 & 0.13 & (0.12, 0.14) & High-middle \\
31 & Estonia & 98.10 & 97.44 & 0.66 & 0.04 & (0.02, 0.07) & High \\
32 & Bulgaria & 98.09 & 96.38 & 1.71 & 0.09 & (0.08, 0.11) & Middle \\
33 & Latvia & 98.03 & 97.77 & 0.26 & 0.03 & (0.01, 0.05) & High-middle \\
34 & Belarus & 98.00 & 97.58 & 0.42 & 0.03 & (0.01, 0.06) & High-middle \\
35 & Turkey & 97.93 & 93.47 & 4.47 & 0.14 & (0.12, 0.17) & High-middle \\
36 & Israel & 97.89 & 95.73 & 2.16 & 0.10 & (0.09, 0.11) & High \\
37 & Lithuania & 97.88 & 97.03 & 0.85 & 0.05 & (0.02, 0.08) & High-middle \\
38 & Ireland & 97.85 & 94.64 & 3.21 & 0.13 & (0.12, 0.14) & High \\
39 & Syria & 97.80 & 95.76 & 2.04 & 0.07 & (0.06, 0.08) & Low-middle \\
40 & Hungary & 97.79 & 93.99 & 3.80 & 0.14 & (0.12, 0.15) & High-middle \\
41 & United States & 97.72 & 94.56 & 3.16 & 0.09 & (0.07, 0.11) & High \\
42 & Thailand & 97.71 & 94.00 & 3.72 & 0.14 & (0.13, 0.15) & Middle \\
43 & Spain & 97.68 & 95.37 & 2.31 & 0.09 & (0.08, 0.10) & High \\
44 & Iceland & 97.68 & 95.73 & 1.94 & 0.09 & (0.07, 0.10) & High \\
45 & Kazakhstan & 97.67 & 95.99 & 1.68 & 0.19 & (0.12, 0.25) & Middle \\
46 & Maldives & 97.63 & 91.36 & 6.28 & 0.18 & (0.16, 0.21) & High-middle \\
47 & Slovakia & 97.63 & 95.08 & 2.55 & 0.07 & (0.07, 0.08) & High \\
48 & Oman & 97.60 & 94.63 & 2.97 & 0.07 & (0.05, 0.08) & High-middle \\
49 & Seychelles & 97.58 & 95.92 & 1.66 & 0.07 & (0.06, 0.07) & High-middle \\
50 & France & 97.55 & 94.77 & 2.78 & 0.13 & (0.12, 0.15) & High \\
51 & Russian Federation & 97.53 & 96.74 & 0.80 & 0.05 & (0.02, 0.07) & High-middle \\
52 & Puerto Rico & 97.51 & 93.43 & 4.08 & 0.14 & (0.11, 0.16) & High \\
53 & Ecuador & 97.51 & 91.59 & 5.92 & 0.18 & (0.16, 0.19) & Middle \\
54 & Azerbaijan & 97.49 & 94.35 & 3.14 & 0.17 & (0.14, 0.19) & Middle \\
55 & Palestine & 97.41 & 95.30 & 2.11 & 0.06 & (0.05, 0.07) & Low-middle \\
56 & Macedonia & 97.40 & 93.54 & 3.85 & 0.12 & (0.11, 0.14) & Middle \\
57 & Serbia & 97.36 & 94.69 & 2.67 & 0.11 & (0.09, 0.12) & Middle \\
58 & Malaysia & 97.33 & 95.06 & 2.26 & 0.08 & (0.07, 0.08) & High-middle \\
59 & Poland & 97.32 & 94.75 & 2.57 & 0.09 & (0.07, 0.11) & High-middle \\
60 & Croatia & 97.26 & 93.54 & 3.72 & 0.10 & (0.09, 0.11) & High-middle \\
61 & Ukraine & 97.25 & 96.71 & 0.54 & 0.05 & (0.02, 0.08) & Middle \\
62 & Kyrgyzstan & 97.22 & 95.61 & 1.61 & 0.14 & (0.09, 0.19) & Low-middle \\
63 & Finland & 97.21 & 93.32 & 3.89 & 0.14 & (0.14, 0.15) & High \\
64 & Virgin Islands, U.S. & 97.20 & 95.78 & 1.41 & 0.05 & (0.04, 0.06) & High \\
65 & American Samoa & 97.19 & 95.96 & 1.24 & 0.06 & (0.05, 0.07) & Middle \\
66 & Philippines & 97.19 & 93.72 & 3.47 & 0.10 & (0.09, 0.12) & Middle \\
67 & Norway & 97.18 & 96.23 & 0.95 & 0.04 & (0.03, 0.05) & High \\
68 & Kuwait & 97.18 & 96.93 & 0.25 & 0.02 & (0.01, 0.04) & High-middle \\
69 & Uzbekistan & 97.14 & 95.36 & 1.79 & 0.11 & (0.07, 0.15) & Middle \\
70 & Jamaica & 97.13 & 95.71 & 1.42 & 0.05 & (0.03, 0.06) & Middle \\
71 & Japan & 97.09 & 95.67 & 1.41 & 0.06 & (0.05, 0.06) & High \\
72 & Portugal & 97.08 & 95.25 & 1.84 & 0.09 & (0.08, 0.10) & High \\
73 & Egypt & 97.07 & 94.23 & 2.84 & 0.08 & (0.07, 0.09) & Middle \\
74 & Trinidad and Tobago & 97.04 & 94.64 & 2.40 & 0.09 & (0.08, 0.11) & High-middle \\
75 & Andorra & 97.03 & 96.54 & 0.50 & 0.03 & (0.01, 0.04) & High \\
76 & Guam & 97.03 & 95.91 & 1.12 & 0.06 & (0.04, 0.07) & High-middle \\
77 & Tonga & 96.89 & 95.60 & 1.29 & 0.03 & (0.03, 0.04) & Middle \\
78 & Uruguay & 96.85 & 94.68 & 2.17 & 0.08 & (0.07, 0.09) & High \\
79 & Colombia & 96.84 & 94.29 & 2.55 & 0.09 & (0.08, 0.11) & Middle \\
80 & Jordan & 96.82 & 94.94 & 1.88 & 0.07 & (0.06, 0.08) & Middle \\
81 & Bahrain & 96.73 & 94.12 & 2.60 & 0.10 & (0.09, 0.10) & High \\
82 & Cyprus & 96.72 & 91.65 & 5.07 & 0.16 & (0.14, 0.18) & High \\
83 & Bosnia and Herzegovina & 96.71 & 93.74 & 2.97 & 0.11 & (0.09, 0.13) & Middle \\
84 & Saint Vincent and the Grenadines & 96.71 & 93.27 & 3.44 & 0.10 & (0.09, 0.12) & Middle \\
85 & Fiji & 96.71 & 95.87 & 0.84 & 0.03 & (0.02, 0.03) & Middle \\
86 & Vietnam & 96.65 & 91.80 & 4.85 & 0.16 & (0.14, 0.17) & Middle \\
87 & Cape Verde & 96.60 & 95.94 & 0.66 & 0.04 & (0.03, 0.05) & Middle \\
88 & South Africa & 96.57 & 95.17 & 1.40 & 0.08 & (0.03, 0.12) & High-middle \\
89 & Greece & 96.57 & 93.83 & 2.74 & 0.09 & (0.07, 0.12) & High \\
90 & Romania & 96.56 & 95.38 & 1.18 & 0.09 & (0.07, 0.12) & Middle \\
91 & Brazil & 96.55 & 93.91 & 2.63 & 0.10 & (0.09, 0.11) & High-middle \\
92 & South Korea & 96.54 & 90.92 & 5.62 & 0.18 & (0.16, 0.20) & High \\
93 & Costa Rica & 96.43 & 93.54 & 2.90 & 0.10 & (0.08, 0.12) & Middle \\
94 & Mexico & 96.38 & 88.08 & 8.30 & 0.24 & (0.19, 0.28) & High-middle \\
95 & Suriname & 96.31 & 94.36 & 1.96 & 0.07 & (0.06, 0.09) & Middle \\
96 & Samoa & 96.28 & 94.61 & 1.67 & 0.05 & (0.04, 0.06) & Middle \\
97 & The Bahamas & 96.18 & 93.66 & 2.52 & 0.11 & (0.10, 0.13) & High-middle \\
98 & Peru & 96.10 & 90.80 & 5.30 & 0.18 & (0.15, 0.22) & Middle \\
99 & Iraq & 96.09 & 94.03 & 2.06 & 0.10 & (0.09, 0.11) & Low-middle \\
100 & Qatar & 96.05 & 91.50 & 4.55 & 0.19 & (0.18, 0.20) & High \\
101 & Federated States of Micronesia & 96.04 & 93.83 & 2.21 & 0.08 & (0.07, 0.08) & Low-middle \\
102 & Zambia & 96.04 & 88.62 & 7.42 & 0.35 & (0.30, 0.40) & Low \\
103 & Saudi Arabia & 96.03 & 93.05 & 2.98 & 0.10 & (0.09, 0.10) & High-middle \\
104 & Panama & 96.02 & 92.01 & 4.01 & 0.14 & (0.12, 0.16) & High-middle \\
105 & Albania & 95.96 & 91.28 & 4.69 & 0.17 & (0.14, 0.19) & Middle \\
106 & Montenegro & 95.94 & 95.02 & 0.92 & 0.05 & (0.03, 0.07) & Middle \\
107 & United Arab Emirates & 95.93 & 95.22 & 0.70 & 0.01 & (-0.00, 0.03) & High \\
108 & Argentina & 95.90 & 92.34 & 3.56 & 0.12 & (0.10, 0.14) & High-middle \\
109 & Saint Lucia & 95.79 & 92.16 & 3.63 & 0.12 & (0.10, 0.14) & Middle \\
110 & Northern Mariana Islands & 95.75 & 92.77 & 2.98 & 0.11 & (0.09, 0.12) & High \\
111 & Tunisia & 95.72 & 93.68 & 2.04 & 0.05 & (0.04, 0.06) & Middle \\
112 & Brunei & 95.70 & 92.67 & 3.02 & 0.07 & (0.05, 0.10) & High-middle \\
113 & Equatorial Guinea & 95.68 & 87.93 & 7.75 & 0.31 & (0.26, 0.36) & Low-middle \\
114 & Sao Tome and Principe & 95.61 & 94.76 & 0.85 & 0.05 & (0.04, 0.06) & Low \\
115 & Cambodia & 95.41 & 92.59 & 2.82 & 0.12 & (0.11, 0.13) & Low-middle \\
116 & Solomon Islands & 95.40 & 93.46 & 1.94 & 0.05 & (0.05, 0.06) & Low \\
117 & Nepal & 95.37 & 92.07 & 3.29 & 0.10 & (0.08, 0.12) & Low-middle \\
118 & Belize & 95.36 & 92.79 & 2.57 & 0.10 & (0.08, 0.11) & Middle \\
119 & Venezuela & 95.33 & 93.16 & 2.17 & 0.08 & (0.06, 0.10) & Middle \\
120 & India & 95.31 & 90.87 & 4.43 & 0.14 & (0.13, 0.14) & Low-middle \\
121 & Turkmenistan & 95.25 & 93.85 & 1.40 & 0.14 & (0.09, 0.19) & Middle \\
122 & Paraguay & 95.23 & 93.11 & 2.11 & 0.06 & (0.06, 0.07) & Middle \\
123 & Bangladesh & 95.21 & 91.35 & 3.86 & 0.13 & (0.12, 0.15) & Low-middle \\
124 & Bhutan & 95.21 & 89.80 & 5.41 & 0.20 & (0.20, 0.21) & Low-middle \\
125 & Myanmar & 95.19 & 89.90 & 5.29 & 0.22 & (0.20, 0.24) & Low-middle \\
126 & The Gambia & 95.08 & 95.42 & -0.33 & 0.01 & (-0.00, 0.02) & Low \\
127 & Dominica & 95.04 & 93.21 & 1.83 & 0.06 & (0.04, 0.07) & Middle \\
128 & Liberia & 95.03 & 93.96 & 1.07 & 0.07 & (0.06, 0.09) & Low \\
129 & Chile & 95.01 & 89.48 & 5.53 & 0.15 & (0.14, 0.17) & High \\
130 & Greenland & 94.96 & 87.42 & 7.54 & 0.29 & (0.25, 0.34) & High \\
131 & Gabon & 94.95 & 91.75 & 3.20 & 0.11 & (0.11, 0.12) & Middle \\
132 & Yemen & 94.94 & 92.32 & 2.62 & 0.10 & (0.08, 0.11) & Low \\
133 & Dominican Republic & 94.84 & 91.95 & 2.89 & 0.07 & (0.05, 0.09) & Middle \\
134 & Botswana & 94.73 & 90.09 & 4.64 & 0.17 & (0.15, 0.19) & Middle \\
135 & Benin & 94.61 & 92.82 & 1.79 & 0.07 & (0.06, 0.08) & Low \\
136 & Mauritania & 94.60 & 92.10 & 2.50 & 0.09 & (0.08, 0.10) & Low-middle \\
137 & Cote d'Ivoire & 94.52 & 93.75 & 0.78 & 0.04 & (0.03, 0.06) & Low-middle \\
138 & Libya & 94.43 & 93.88 & 0.55 & -0.01 & (-0.03, 0.01) & Middle \\
139 & Namibia & 94.42 & 93.89 & 0.52 & 0.08 & (0.03, 0.12) & Low-middle \\
140 & Vanuatu & 94.41 & 92.95 & 1.46 & 0.05 & (0.04, 0.05) & Low \\
141 & Indonesia & 94.40 & 91.79 & 2.61 & 0.06 & (0.05, 0.07) & Middle \\
142 & Sierra Leone & 94.30 & 93.57 & 0.73 & 0.02 & (0.01, 0.02) & Low \\
143 & Sudan & 94.25 & 88.74 & 5.51 & 0.17 & (0.15, 0.19) & Low \\
144 & Senegal & 94.22 & 92.88 & 1.35 & 0.04 & (0.02, 0.05) & Low-middle \\
145 & Cameroon & 94.21 & 92.50 & 1.71 & 0.06 & (0.05, 0.07) & Low-middle \\
146 & Nicaragua & 94.16 & 88.49 & 5.68 & 0.20 & (0.17, 0.22) & Low-middle \\
147 & Guyana & 94.15 & 89.04 & 5.10 & 0.15 & (0.12, 0.18) & Low-middle \\
148 & Congo & 94.13 & 90.33 & 3.80 & 0.17 & (0.15, 0.19) & Low \\
149 & Togo & 94.05 & 93.91 & 0.15 & 0.01 & (0.01, 0.02) & Low \\
150 & Angola & 93.94 & 88.75 & 5.19 & 0.20 & (0.19, 0.21) & Low \\
151 & Djibouti & 93.89 & 93.02 & 0.86 & 0.05 & (0.03, 0.07) & Low \\
152 & Guatemala & 93.87 & 81.21 & 12.66 & 0.47 & (0.36, 0.58) & Low-middle \\
153 & Ethiopia & 93.85 & 85.55 & 8.30 & 0.34 & (0.32, 0.36) & Low \\
154 & Ghana & 93.81 & 92.48 & 1.33 & 0.04 & (0.03, 0.04) & Low-middle \\
155 & Papua New Guinea & 93.75 & 90.90 & 2.85 & 0.09 & (0.08, 0.10) & Low \\
156 & Eritrea & 93.73 & 92.55 & 1.18 & 0.07 & (0.06, 0.08) & Low \\
157 & Comoros & 93.63 & 91.02 & 2.61 & 0.12 & (0.10, 0.13) & Low-middle \\
158 & El Salvador & 93.62 & 86.72 & 6.90 & 0.22 & (0.17, 0.28) & Middle \\
159 & Democratic Republic of the Congo & 93.62 & 91.04 & 2.58 & 0.09 & (0.08, 0.11) & Low \\
160 & Guinea & 93.49 & 91.63 & 1.86 & 0.08 & (0.07, 0.08) & Low \\
161 & Madagascar & 93.35 & 92.25 & 1.10 & 0.05 & (0.04, 0.05) & Low \\
162 & Bolivia & 93.34 & 85.70 & 7.63 & 0.26 & (0.22, 0.29) & Low-middle \\
163 & Mali & 93.28 & 91.62 & 1.66 & 0.07 & (0.06, 0.08) & Low \\
164 & Burkina Faso & 93.28 & 93.18 & 0.09 & -0.01 & (-0.03, 0.01) & Low \\
165 & Niger & 93.25 & 90.89 & 2.36 & 0.11 & (0.09, 0.13) & Low \\
166 & Nigeria & 93.17 & 91.27 & 1.90 & 0.09 & (0.08, 0.10) & Low \\
167 & Laos & 93.14 & 89.07 & 4.07 & 0.14 & (0.13, 0.15) & Low-middle \\
168 & Morocco & 92.96 & 88.05 & 4.92 & 0.15 & (0.13, 0.16) & Middle \\
169 & Tanzania & 92.87 & 91.99 & 0.87 & 0.03 & (0.02, 0.04) & Low-middle \\
170 & North Korea & 92.82 & 89.09 & 3.73 & 0.14 & (0.13, 0.15) & Middle \\
171 & Tajikistan & 92.79 & 89.97 & 2.82 & 0.20 & (0.16, 0.24) & Low-middle \\
172 & Pakistan & 92.73 & 91.72 & 1.01 & 0.05 & (0.03, 0.06) & Low-middle \\
173 & Timor-Leste & 92.73 & 87.54 & 5.19 & 0.16 & (0.14, 0.19) & Low \\
174 & Uganda & 92.47 & 89.47 & 3.00 & 0.15 & (0.12, 0.18) & Low \\
175 & Honduras & 92.42 & 87.73 & 4.69 & 0.15 & (0.13, 0.17) & Low-middle \\
176 & South Sudan & 92.36 & 91.88 & 0.48 & 0.04 & (0.02, 0.06) & Low \\
177 & Rwanda & 92.29 & 85.94 & 6.35 & 0.43 & (0.34, 0.53) & Low \\
178 & Burundi & 91.95 & 88.89 & 3.06 & 0.16 & (0.14, 0.18) & Low \\
179 & Algeria & 91.71 & 86.76 & 4.95 & 0.16 & (0.14, 0.18) & High-middle \\
180 & Mongolia & 91.70 & 86.27 & 5.43 & 0.30 & (0.27, 0.34) & Middle \\
181 & Chad & 91.50 & 92.01 & -0.51 & 0.02 & (-0.01, 0.05) & Low \\
182 & Malawi & 91.42 & 89.69 & 1.73 & 0.05 & (0.04, 0.07) & Low \\
183 & Eswatini & 91.30 & 94.30 & -2.99 & -0.07 & (-0.18, 0.04) & Middle \\
184 & Kiribati & 91.01 & 87.35 & 3.66 & 0.12 & (0.11, 0.13) & Low \\
185 & Mozambique & 90.96 & 88.88 & 2.08 & 0.05 & (0.03, 0.07) & Low \\
186 & Haiti & 89.75 & 81.77 & 7.98 & 0.27 & (0.24, 0.31) & Low \\
187 & Kenya & 88.60 & 89.59 & -0.99 & -0.07 & (-0.09, -0.06) & Low-middle \\
188 & Guinea-Bissau & 88.32 & 87.87 & 0.46 & 0.04 & (0.02, 0.07) & Low \\
189 & Somalia & 88.29 & 88.17 & 0.11 & 0.04 & (0.03, 0.05) & Low \\
190 & Central African Republic & 87.17 & 85.61 & 1.56 & 0.11 & (0.09, 0.14) & Low \\
191 & Marshall Islands & 87.16 & 82.48 & 4.67 & 0.16 & (0.14, 0.17) & Low-middle \\
192 & Zimbabwe & 85.51 & 93.33 & -7.81 & -0.36 & (-0.42, -0.30) & Low \\
193 & Lesotho & 84.60 & 91.21 & -6.61 & -0.26 & (-0.34, -0.19) & Low \\
194 & Afghanistan & 82.39 & 76.16 & 6.23 & 0.33 & (0.29, 0.37) & Low \\

\end{longtable}


\begin{longtable}{r l r r r r l l}
\caption{Complete country ranking by Quality of Care Index in 2021 (age-standardised, both sexes), with QCI in 1990, absolute change, Average Annual Percentage Change (AAPC) over 1990--2021, and SDI group.}
\label{tab:all_countries} \\
\toprule
\textbf{Rank} & \textbf{Country} & \textbf{QCI 2021} & \textbf{QCI 1990} & \textbf{Change} & \textbf{AAPC} & \textbf{95\% CI} & \textbf{SDI Group} \\
\midrule
\endfirsthead

\multicolumn{8}{l}{\textit{Table~S3 continued from previous page}} \\
\toprule
\textbf{Rank} & \textbf{Country} & \textbf{QCI 2021} & \textbf{QCI 1990} & \textbf{Change} & \textbf{AAPC} & \textbf{95\% CI} & \textbf{SDI Group} \\
\midrule
\endhead

\midrule
\multicolumn{8}{r}{\textit{Continued on next page}} \\
\endfoot

\bottomrule
\endlastfoot
1 & Bermuda & 99.94 & 99.65 & 0.29 & 0.01 & (0.01, 0.01) & High \\
2 & New Zealand & 99.63 & 97.38 & 2.25 & 0.09 & (0.07, 0.10) & High \\
3 & Malta & 99.54 & 98.60 & 0.94 & 0.04 & (0.03, 0.04) & High \\
4 & Antigua and Barbuda & 99.31 & 98.02 & 1.28 & 0.04 & (0.04, 0.05) & High-middle \\
5 & Luxembourg & 99.16 & 97.62 & 1.54 & 0.06 & (0.06, 0.07) & High \\
6 & Cuba & 99.04 & 97.89 & 1.15 & 0.04 & (0.03, 0.05) & High-middle \\
7 & Taiwan & 99.01 & 93.34 & 5.67 & 0.13 & (0.11, 0.16) & High \\
8 & Sweden & 98.99 & 95.87 & 3.12 & 0.13 & (0.12, 0.14) & High \\
9 & Switzerland & 98.94 & 96.68 & 2.26 & 0.09 & (0.08, 0.09) & High \\
10 & Grenada & 98.91 & 96.49 & 2.42 & 0.07 & (0.06, 0.08) & Middle \\
11 & Czech Republic & 98.79 & 96.86 & 1.92 & 0.06 & (0.05, 0.06) & High \\
12 & Slovenia & 98.70 & 96.64 & 2.07 & 0.07 & (0.06, 0.08) & High \\
13 & Denmark & 98.70 & 97.09 & 1.62 & 0.06 & (0.06, 0.07) & High \\
14 & Australia & 98.70 & 97.36 & 1.34 & 0.05 & (0.05, 0.06) & High \\
15 & Italy & 98.69 & 96.69 & 1.99 & 0.07 & (0.05, 0.08) & High \\
16 & China & 98.68 & 94.21 & 4.47 & 0.14 & (0.12, 0.15) & High-middle \\
17 & Mauritius & 98.64 & 96.71 & 1.94 & 0.05 & (0.04, 0.06) & High-middle \\
18 & Singapore & 98.61 & 95.64 & 2.97 & 0.10 & (0.09, 0.10) & High \\
19 & Armenia & 98.60 & 97.09 & 1.51 & 0.09 & (0.07, 0.11) & Middle \\
20 & Belgium & 98.58 & 96.62 & 1.96 & 0.08 & (0.07, 0.08) & High \\
21 & Germany & 98.49 & 95.65 & 2.84 & 0.10 & (0.09, 0.11) & High \\
22 & Moldova & 98.45 & 97.52 & 0.92 & 0.06 & (0.04, 0.09) & Middle \\
23 & United Kingdom & 98.44 & 97.42 & 1.02 & 0.05 & (0.04, 0.05) & High \\
24 & Canada & 98.35 & 95.04 & 3.31 & 0.11 & (0.09, 0.13) & High \\
25 & Netherlands & 98.28 & 96.57 & 1.70 & 0.08 & (0.07, 0.09) & High \\
26 & Georgia & 98.26 & 95.34 & 2.92 & 0.12 & (0.10, 0.14) & Middle \\
27 & Lebanon & 98.24 & 95.90 & 2.34 & 0.08 & (0.07, 0.08) & Middle \\
28 & Austria & 98.19 & 95.66 & 2.52 & 0.08 & (0.07, 0.10) & High \\
29 & Barbados & 98.16 & 96.78 & 1.38 & 0.06 & (0.05, 0.07) & High \\
30 & Sri Lanka & 98.13 & 93.59 & 4.54 & 0.13 & (0.12, 0.14) & High-middle \\
31 & Estonia & 98.10 & 97.44 & 0.66 & 0.04 & (0.02, 0.07) & High \\
32 & Bulgaria & 98.09 & 96.38 & 1.71 & 0.09 & (0.08, 0.11) & Middle \\
33 & Latvia & 98.03 & 97.77 & 0.26 & 0.03 & (0.01, 0.05) & High-middle \\
34 & Belarus & 98.00 & 97.58 & 0.42 & 0.03 & (0.01, 0.06) & High-middle \\
35 & Turkey & 97.93 & 93.47 & 4.47 & 0.14 & (0.12, 0.17) & High-middle \\
36 & Israel & 97.89 & 95.73 & 2.16 & 0.10 & (0.09, 0.11) & High \\
37 & Lithuania & 97.88 & 97.03 & 0.85 & 0.05 & (0.02, 0.08) & High-middle \\
38 & Ireland & 97.85 & 94.64 & 3.21 & 0.13 & (0.12, 0.14) & High \\
39 & Syria & 97.80 & 95.76 & 2.04 & 0.07 & (0.06, 0.08) & Low-middle \\
40 & Hungary & 97.79 & 93.99 & 3.80 & 0.14 & (0.12, 0.15) & High-middle \\
41 & United States & 97.72 & 94.56 & 3.16 & 0.09 & (0.07, 0.11) & High \\
42 & Thailand & 97.71 & 94.00 & 3.72 & 0.14 & (0.13, 0.15) & Middle \\
43 & Spain & 97.68 & 95.37 & 2.31 & 0.09 & (0.08, 0.10) & High \\
44 & Iceland & 97.68 & 95.73 & 1.94 & 0.09 & (0.07, 0.10) & High \\
45 & Kazakhstan & 97.67 & 95.99 & 1.68 & 0.19 & (0.12, 0.25) & Middle \\
46 & Maldives & 97.63 & 91.36 & 6.28 & 0.18 & (0.16, 0.21) & High-middle \\
47 & Slovakia & 97.63 & 95.08 & 2.55 & 0.07 & (0.07, 0.08) & High \\
48 & Oman & 97.60 & 94.63 & 2.97 & 0.07 & (0.05, 0.08) & High-middle \\
49 & Seychelles & 97.58 & 95.92 & 1.66 & 0.07 & (0.06, 0.07) & High-middle \\
50 & France & 97.55 & 94.77 & 2.78 & 0.13 & (0.12, 0.15) & High \\
51 & Russian Federation & 97.53 & 96.74 & 0.80 & 0.05 & (0.02, 0.07) & High-middle \\
52 & Puerto Rico & 97.51 & 93.43 & 4.08 & 0.14 & (0.11, 0.16) & High \\
53 & Ecuador & 97.51 & 91.59 & 5.92 & 0.18 & (0.16, 0.19) & Middle \\
54 & Azerbaijan & 97.49 & 94.35 & 3.14 & 0.17 & (0.14, 0.19) & Middle \\
55 & Palestine & 97.41 & 95.30 & 2.11 & 0.06 & (0.05, 0.07) & Low-middle \\
56 & Macedonia & 97.40 & 93.54 & 3.85 & 0.12 & (0.11, 0.14) & Middle \\
57 & Serbia & 97.36 & 94.69 & 2.67 & 0.11 & (0.09, 0.12) & Middle \\
58 & Malaysia & 97.33 & 95.06 & 2.26 & 0.08 & (0.07, 0.08) & High-middle \\
59 & Poland & 97.32 & 94.75 & 2.57 & 0.09 & (0.07, 0.11) & High-middle \\
60 & Croatia & 97.26 & 93.54 & 3.72 & 0.10 & (0.09, 0.11) & High-middle \\
61 & Ukraine & 97.25 & 96.71 & 0.54 & 0.05 & (0.02, 0.08) & Middle \\
62 & Kyrgyzstan & 97.22 & 95.61 & 1.61 & 0.14 & (0.09, 0.19) & Low-middle \\
63 & Finland & 97.21 & 93.32 & 3.89 & 0.14 & (0.14, 0.15) & High \\
64 & Virgin Islands, U.S. & 97.20 & 95.78 & 1.41 & 0.05 & (0.04, 0.06) & High \\
65 & American Samoa & 97.19 & 95.96 & 1.24 & 0.06 & (0.05, 0.07) & Middle \\
66 & Philippines & 97.19 & 93.72 & 3.47 & 0.10 & (0.09, 0.12) & Middle \\
67 & Norway & 97.18 & 96.23 & 0.95 & 0.04 & (0.03, 0.05) & High \\
68 & Kuwait & 97.18 & 96.93 & 0.25 & 0.02 & (0.01, 0.04) & High-middle \\
69 & Uzbekistan & 97.14 & 95.36 & 1.79 & 0.11 & (0.07, 0.15) & Middle \\
70 & Jamaica & 97.13 & 95.71 & 1.42 & 0.05 & (0.03, 0.06) & Middle \\
71 & Japan & 97.09 & 95.67 & 1.41 & 0.06 & (0.05, 0.06) & High \\
72 & Portugal & 97.08 & 95.25 & 1.84 & 0.09 & (0.08, 0.10) & High \\
73 & Egypt & 97.07 & 94.23 & 2.84 & 0.08 & (0.07, 0.09) & Middle \\
74 & Trinidad and Tobago & 97.04 & 94.64 & 2.40 & 0.09 & (0.08, 0.11) & High-middle \\
75 & Andorra & 97.03 & 96.54 & 0.50 & 0.03 & (0.01, 0.04) & High \\
76 & Guam & 97.03 & 95.91 & 1.12 & 0.06 & (0.04, 0.07) & High-middle \\
77 & Tonga & 96.89 & 95.60 & 1.29 & 0.03 & (0.03, 0.04) & Middle \\
78 & Uruguay & 96.85 & 94.68 & 2.17 & 0.08 & (0.07, 0.09) & High \\
79 & Colombia & 96.84 & 94.29 & 2.55 & 0.09 & (0.08, 0.11) & Middle \\
80 & Jordan & 96.82 & 94.94 & 1.88 & 0.07 & (0.06, 0.08) & Middle \\
81 & Bahrain & 96.73 & 94.12 & 2.60 & 0.10 & (0.09, 0.10) & High \\
82 & Cyprus & 96.72 & 91.65 & 5.07 & 0.16 & (0.14, 0.18) & High \\
83 & Bosnia and Herzegovina & 96.71 & 93.74 & 2.97 & 0.11 & (0.09, 0.13) & Middle \\
84 & Saint Vincent and the Grenadines & 96.71 & 93.27 & 3.44 & 0.10 & (0.09, 0.12) & Middle \\
85 & Fiji & 96.71 & 95.87 & 0.84 & 0.03 & (0.02, 0.03) & Middle \\
86 & Vietnam & 96.65 & 91.80 & 4.85 & 0.16 & (0.14, 0.17) & Middle \\
87 & Cape Verde & 96.60 & 95.94 & 0.66 & 0.04 & (0.03, 0.05) & Middle \\
88 & South Africa & 96.57 & 95.17 & 1.40 & 0.08 & (0.03, 0.12) & High-middle \\
89 & Greece & 96.57 & 93.83 & 2.74 & 0.09 & (0.07, 0.12) & High \\
90 & Romania & 96.56 & 95.38 & 1.18 & 0.09 & (0.07, 0.12) & Middle \\
91 & Brazil & 96.55 & 93.91 & 2.63 & 0.10 & (0.09, 0.11) & High-middle \\
92 & South Korea & 96.54 & 90.92 & 5.62 & 0.18 & (0.16, 0.20) & High \\
93 & Costa Rica & 96.43 & 93.54 & 2.90 & 0.10 & (0.08, 0.12) & Middle \\
94 & Mexico & 96.38 & 88.08 & 8.30 & 0.24 & (0.19, 0.28) & High-middle \\
95 & Suriname & 96.31 & 94.36 & 1.96 & 0.07 & (0.06, 0.09) & Middle \\
96 & Samoa & 96.28 & 94.61 & 1.67 & 0.05 & (0.04, 0.06) & Middle \\
97 & The Bahamas & 96.18 & 93.66 & 2.52 & 0.11 & (0.10, 0.13) & High-middle \\
98 & Peru & 96.10 & 90.80 & 5.30 & 0.18 & (0.15, 0.22) & Middle \\
99 & Iraq & 96.09 & 94.03 & 2.06 & 0.10 & (0.09, 0.11) & Low-middle \\
100 & Qatar & 96.05 & 91.50 & 4.55 & 0.19 & (0.18, 0.20) & High \\
101 & Federated States of Micronesia & 96.04 & 93.83 & 2.21 & 0.08 & (0.07, 0.08) & Low-middle \\
102 & Zambia & 96.04 & 88.62 & 7.42 & 0.35 & (0.30, 0.40) & Low \\
103 & Saudi Arabia & 96.03 & 93.05 & 2.98 & 0.10 & (0.09, 0.10) & High-middle \\
104 & Panama & 96.02 & 92.01 & 4.01 & 0.14 & (0.12, 0.16) & High-middle \\
105 & Albania & 95.96 & 91.28 & 4.69 & 0.17 & (0.14, 0.19) & Middle \\
106 & Montenegro & 95.94 & 95.02 & 0.92 & 0.05 & (0.03, 0.07) & Middle \\
107 & United Arab Emirates & 95.93 & 95.22 & 0.70 & 0.01 & (-0.00, 0.03) & High \\
108 & Argentina & 95.90 & 92.34 & 3.56 & 0.12 & (0.10, 0.14) & High-middle \\
109 & Saint Lucia & 95.79 & 92.16 & 3.63 & 0.12 & (0.10, 0.14) & Middle \\
110 & Northern Mariana Islands & 95.75 & 92.77 & 2.98 & 0.11 & (0.09, 0.12) & High \\
111 & Tunisia & 95.72 & 93.68 & 2.04 & 0.05 & (0.04, 0.06) & Middle \\
112 & Brunei & 95.70 & 92.67 & 3.02 & 0.07 & (0.05, 0.10) & High-middle \\
113 & Equatorial Guinea & 95.68 & 87.93 & 7.75 & 0.31 & (0.26, 0.36) & Low-middle \\
114 & Sao Tome and Principe & 95.61 & 94.76 & 0.85 & 0.05 & (0.04, 0.06) & Low \\
115 & Cambodia & 95.41 & 92.59 & 2.82 & 0.12 & (0.11, 0.13) & Low-middle \\
116 & Solomon Islands & 95.40 & 93.46 & 1.94 & 0.05 & (0.05, 0.06) & Low \\
117 & Nepal & 95.37 & 92.07 & 3.29 & 0.10 & (0.08, 0.12) & Low-middle \\
118 & Belize & 95.36 & 92.79 & 2.57 & 0.10 & (0.08, 0.11) & Middle \\
119 & Venezuela & 95.33 & 93.16 & 2.17 & 0.08 & (0.06, 0.10) & Middle \\
120 & India & 95.31 & 90.87 & 4.43 & 0.14 & (0.13, 0.14) & Low-middle \\
121 & Turkmenistan & 95.25 & 93.85 & 1.40 & 0.14 & (0.09, 0.19) & Middle \\
122 & Paraguay & 95.23 & 93.11 & 2.11 & 0.06 & (0.06, 0.07) & Middle \\
123 & Bangladesh & 95.21 & 91.35 & 3.86 & 0.13 & (0.12, 0.15) & Low-middle \\
124 & Bhutan & 95.21 & 89.80 & 5.41 & 0.20 & (0.20, 0.21) & Low-middle \\
125 & Myanmar & 95.19 & 89.90 & 5.29 & 0.22 & (0.20, 0.24) & Low-middle \\
126 & The Gambia & 95.08 & 95.42 & -0.33 & 0.01 & (-0.00, 0.02) & Low \\
127 & Dominica & 95.04 & 93.21 & 1.83 & 0.06 & (0.04, 0.07) & Middle \\
128 & Liberia & 95.03 & 93.96 & 1.07 & 0.07 & (0.06, 0.09) & Low \\
129 & Chile & 95.01 & 89.48 & 5.53 & 0.15 & (0.14, 0.17) & High \\
130 & Greenland & 94.96 & 87.42 & 7.54 & 0.29 & (0.25, 0.34) & High \\
131 & Gabon & 94.95 & 91.75 & 3.20 & 0.11 & (0.11, 0.12) & Middle \\
132 & Yemen & 94.94 & 92.32 & 2.62 & 0.10 & (0.08, 0.11) & Low \\
133 & Dominican Republic & 94.84 & 91.95 & 2.89 & 0.07 & (0.05, 0.09) & Middle \\
134 & Botswana & 94.73 & 90.09 & 4.64 & 0.17 & (0.15, 0.19) & Middle \\
135 & Benin & 94.61 & 92.82 & 1.79 & 0.07 & (0.06, 0.08) & Low \\
136 & Mauritania & 94.60 & 92.10 & 2.50 & 0.09 & (0.08, 0.10) & Low-middle \\
137 & Cote d'Ivoire & 94.52 & 93.75 & 0.78 & 0.04 & (0.03, 0.06) & Low-middle \\
138 & Libya & 94.43 & 93.88 & 0.55 & -0.01 & (-0.03, 0.01) & Middle \\
139 & Namibia & 94.42 & 93.89 & 0.52 & 0.08 & (0.03, 0.12) & Low-middle \\
140 & Vanuatu & 94.41 & 92.95 & 1.46 & 0.05 & (0.04, 0.05) & Low \\
141 & Indonesia & 94.40 & 91.79 & 2.61 & 0.06 & (0.05, 0.07) & Middle \\
142 & Sierra Leone & 94.30 & 93.57 & 0.73 & 0.02 & (0.01, 0.02) & Low \\
143 & Sudan & 94.25 & 88.74 & 5.51 & 0.17 & (0.15, 0.19) & Low \\
144 & Senegal & 94.22 & 92.88 & 1.35 & 0.04 & (0.02, 0.05) & Low-middle \\
145 & Cameroon & 94.21 & 92.50 & 1.71 & 0.06 & (0.05, 0.07) & Low-middle \\
146 & Nicaragua & 94.16 & 88.49 & 5.68 & 0.20 & (0.17, 0.22) & Low-middle \\
147 & Guyana & 94.15 & 89.04 & 5.10 & 0.15 & (0.12, 0.18) & Low-middle \\
148 & Congo & 94.13 & 90.33 & 3.80 & 0.17 & (0.15, 0.19) & Low \\
149 & Togo & 94.05 & 93.91 & 0.15 & 0.01 & (0.01, 0.02) & Low \\
150 & Angola & 93.94 & 88.75 & 5.19 & 0.20 & (0.19, 0.21) & Low \\
151 & Djibouti & 93.89 & 93.02 & 0.86 & 0.05 & (0.03, 0.07) & Low \\
152 & Guatemala & 93.87 & 81.21 & 12.66 & 0.47 & (0.36, 0.58) & Low-middle \\
153 & Ethiopia & 93.85 & 85.55 & 8.30 & 0.34 & (0.32, 0.36) & Low \\
154 & Ghana & 93.81 & 92.48 & 1.33 & 0.04 & (0.03, 0.04) & Low-middle \\
155 & Papua New Guinea & 93.75 & 90.90 & 2.85 & 0.09 & (0.08, 0.10) & Low \\
156 & Eritrea & 93.73 & 92.55 & 1.18 & 0.07 & (0.06, 0.08) & Low \\
157 & Comoros & 93.63 & 91.02 & 2.61 & 0.12 & (0.10, 0.13) & Low-middle \\
158 & El Salvador & 93.62 & 86.72 & 6.90 & 0.22 & (0.17, 0.28) & Middle \\
159 & Democratic Republic of the Congo & 93.62 & 91.04 & 2.58 & 0.09 & (0.08, 0.11) & Low \\
160 & Guinea & 93.49 & 91.63 & 1.86 & 0.08 & (0.07, 0.08) & Low \\
161 & Madagascar & 93.35 & 92.25 & 1.10 & 0.05 & (0.04, 0.05) & Low \\
162 & Bolivia & 93.34 & 85.70 & 7.63 & 0.26 & (0.22, 0.29) & Low-middle \\
163 & Mali & 93.28 & 91.62 & 1.66 & 0.07 & (0.06, 0.08) & Low \\
164 & Burkina Faso & 93.28 & 93.18 & 0.09 & -0.01 & (-0.03, 0.01) & Low \\
165 & Niger & 93.25 & 90.89 & 2.36 & 0.11 & (0.09, 0.13) & Low \\
166 & Nigeria & 93.17 & 91.27 & 1.90 & 0.09 & (0.08, 0.10) & Low \\
167 & Laos & 93.14 & 89.07 & 4.07 & 0.14 & (0.13, 0.15) & Low-middle \\
168 & Morocco & 92.96 & 88.05 & 4.92 & 0.15 & (0.13, 0.16) & Middle \\
169 & Tanzania & 92.87 & 91.99 & 0.87 & 0.03 & (0.02, 0.04) & Low-middle \\
170 & North Korea & 92.82 & 89.09 & 3.73 & 0.14 & (0.13, 0.15) & Middle \\
171 & Tajikistan & 92.79 & 89.97 & 2.82 & 0.20 & (0.16, 0.24) & Low-middle \\
172 & Pakistan & 92.73 & 91.72 & 1.01 & 0.05 & (0.03, 0.06) & Low-middle \\
173 & Timor-Leste & 92.73 & 87.54 & 5.19 & 0.16 & (0.14, 0.19) & Low \\
174 & Uganda & 92.47 & 89.47 & 3.00 & 0.15 & (0.12, 0.18) & Low \\
175 & Honduras & 92.42 & 87.73 & 4.69 & 0.15 & (0.13, 0.17) & Low-middle \\
176 & South Sudan & 92.36 & 91.88 & 0.48 & 0.04 & (0.02, 0.06) & Low \\
177 & Rwanda & 92.29 & 85.94 & 6.35 & 0.43 & (0.34, 0.53) & Low \\
178 & Burundi & 91.95 & 88.89 & 3.06 & 0.16 & (0.14, 0.18) & Low \\
179 & Algeria & 91.71 & 86.76 & 4.95 & 0.16 & (0.14, 0.18) & High-middle \\
180 & Mongolia & 91.70 & 86.27 & 5.43 & 0.30 & (0.27, 0.34) & Middle \\
181 & Chad & 91.50 & 92.01 & -0.51 & 0.02 & (-0.01, 0.05) & Low \\
182 & Malawi & 91.42 & 89.69 & 1.73 & 0.05 & (0.04, 0.07) & Low \\
183 & Eswatini & 91.30 & 94.30 & -2.99 & -0.07 & (-0.18, 0.04) & Middle \\
184 & Kiribati & 91.01 & 87.35 & 3.66 & 0.12 & (0.11, 0.13) & Low \\
185 & Mozambique & 90.96 & 88.88 & 2.08 & 0.05 & (0.03, 0.07) & Low \\
186 & Haiti & 89.75 & 81.77 & 7.98 & 0.27 & (0.24, 0.31) & Low \\
187 & Kenya & 88.60 & 89.59 & -0.99 & -0.07 & (-0.09, -0.06) & Low-middle \\
188 & Guinea-Bissau & 88.32 & 87.87 & 0.46 & 0.04 & (0.02, 0.07) & Low \\
189 & Somalia & 88.29 & 88.17 & 0.11 & 0.04 & (0.03, 0.05) & Low \\
190 & Central African Republic & 87.17 & 85.61 & 1.56 & 0.11 & (0.09, 0.14) & Low \\
191 & Marshall Islands & 87.16 & 82.48 & 4.67 & 0.16 & (0.14, 0.17) & Low-middle \\
192 & Zimbabwe & 85.51 & 93.33 & -7.81 & -0.36 & (-0.42, -0.30) & Low \\
193 & Lesotho & 84.60 & 91.21 & -6.61 & -0.26 & (-0.34, -0.19) & Low \\
194 & Afghanistan & 82.39 & 76.16 & 6.23 & 0.33 & (0.29, 0.37) & Low \\

\end{longtable}


\end{document}
